% ==============================================================================
% SYSTEM AI PROMPT & TYPESETTING RULES (STRICTLY ENFORCED)
% ==============================================================================
% ИНСТРУКЦИЯ ДЛЯ ИИ-ГЕНЕРАТОРА КОДА. ЧИТАТЬ ОБЯЗАТЕЛЬНО!
% Вы — профессиональный LaTeX-верстальщик. При генерации контента внутри 
% этого документа вы ОБЯЗАНЫ соблюдать следующие правила:
%
% 1. [КРИТИЧЕСКИ ВАЖНО] ПОЛНЫЙ ЗАПРЕТ НА MARKDOWN:
%    - ЗАПРЕЩЕНО использовать `**жирный текст**` (используй \textbf{...})
%    - ЗАПРЕЩЕНО использовать `*курсив*` (используй \textit{...} или \emph{...})
%    - ЗАПРЕЩЕНО использовать `# Заголовок` (используй \section{...}, \subsection{...})
%    - ЗАПРЕЩЕНО использовать обратные кавычки `код` (используй \texttt{...})
%    Любое использование Markdown будет считаться фатальной ошибкой генерации.
%
% 2. СТРОГИЕ СПИСКИ:
%    - Никаких списков через дефисы "-" или звездочки "* " текстом.
%    - Все списки ДОЛЖНЫ быть обернуты в окружения \begin{itemize} ... \end{itemize} 
%      или \begin{enumerate} ... \end{enumerate}. Элементы задаются через \item.
%
% 3. МАТЕМАТИЧЕСКАЯ СЕМАНТИКА И ОКРУЖЕНИЯ:
%    - Используй \begin{equation}, \begin{align}, \begin{gather} для формул.
%    - ЗАПРЕЩЕНО использовать устаревший синтаксис $$...$$ для выключных формул.
%      Используй \[ ... \] или окружение equation*.
%    - Системы уравнений и кусочные функции оформлять ТОЛЬКО через \begin{cases}.
%    - Матрицы оформлять через pmatrix, bmatrix, vmatrix или pNiceMatrix (для блочных).
%
% 4. МИКРОТИПОГРАФИКА И ИНТЕРВАЛЫ:
%    - Для контроля отступов использовать \vspace{...} и \hspace{...}.
%    - Не использовать пустые строки для создания вертикального отступа.
%    - После формул корректно расставлять знаки препинания (запятые, точки).
%
% 5. СТРУКТУРНЫЕ БЛОКИ (ТЕОРЕМЫ И ОПРЕДЕЛЕНИЯ):
%    - Вся теория должна оформляться через семантические окружения:
%      \begin{definition}, \begin{theorem}, \begin{lemma}, \begin{proof}, \begin{remark}.
%    - Для выделенных блоков ОБЯЗАТЕЛЬНО использовать tcolorbox-окружения:
%      \begin{tcbdefinition}, \begin{tcbtheorem}, \begin{tcblemma},
%      \begin{tcbcorollary}, \begin{tcbalgorithmbox}, \begin{tcbwarning},
%      \begin{tcbexample}, \begin{tcbinsight}, \begin{tcbproof}.
%
% 6. [КРИТИЧЕСКИ ВАЖНО] АЛГОРИТМЫ И ПОСЛЕДОВАТЕЛЬНЫЕ ТРЕБОВАНИЯ:
%    - ЛЮБОЙ алгоритм, псевдокод, пошаговая инструкция или нумерованный список
%      последовательных шагов/требований ОБЯЗАН быть обёрнут в \begin{tcbalgorithmbox}.
%    - ЗАПРЕЩЕНО выносить алгоритмы, шаги и последовательности в голый текст или
%      обычный \begin{enumerate} без обёртки в tcbalgorithmbox.
%    - Если перечисление является логически упорядоченной процедурой (даже из 2 шагов),
%      оно считается алгоритмом и требует tcbalgorithmbox.
%
% 7. [КРИТИЧЕСКИ ВАЖНО] ДОКАЗАТЕЛЬСТВА:
%    - Все доказательства оформляются исключительно через \begin{tcbproof} ... \end{tcbproof}.
%    - ЗАПРЕЩЕНО использовать стандартное amsthm окружение \begin{proof}.
%    - Знак конца доказательства \qed расставляется АВТОМАТИЧЕСКИ окружением tcbproof,
%      вручную его добавлять ЗАПРЕЩЕНО.
%
% 8. [КРИТИЧЕСКИ ВАЖНО] УМНЫЕ ССЫЛКИ (\cref и \Cref):
%    - Для ссылок на теоремы, леммы, определения и алгоритмы ВСЕГДА использовать
%      \cref{label} (строчная) или \Cref{label} (в начале предложения).
%    - ЗАПРЕЩЕНО использовать \ref без \cref — это приводит к потере семантики.
%    - Каждый tcolorbox-блок, на который планируется ссылаться, ДОЛЖЕН иметь
%      метку: добавляй \label{...} СРАЗУ после открывающего тега окружения.
%    - Соглашение об именовании меток:
%        определения  — \label{def:...}
%        теоремы      — \label{thm:...}
%        леммы        — \label{lem:...}
%        следствия    — \label{cor:...}
%        алгоритмы    — \label{alg:...}
%        примеры      — \label{ex:...}
%        уравнения    — \label{eq:...}
%        разделы      — \label{sec:...}
%    - ЗАПРЕЩЕНО использовать произвольные префиксы вне этого соглашения.
%
% 9. СПЕЦИФИКА МАТРИЧНЫХ ВЫЧИСЛЕНИЙ (ИСПОЛЬЗУЙ МАКРОСЫ):
%    - Транспонирование: A\T (не A^T и не A^\top).
%    - Эрмитово сопряжение: A\H.
%    - Норма: \norm{A}, модуль: \abs{x}, скалярное произведение: \inner{x, y}.
%    - Операторы: \rank(A), \tr(A), \diag(x), \im(A), \ker(A).
% ==============================================================================

\documentclass[12pt, a4paper]{article}

% --- Кодировки, шрифты и языки ---
\usepackage[utf8]{inputenc}
\usepackage[T2A]{fontenc}
\usepackage[russian, english]{babel}
\usepackage{cmap}
\usepackage{lmodern}

% --- Микротипографика ---
\usepackage{microtype}
\usepackage{indentfirst}

% ==============================================================================
% УЛУЧШЕННАЯ ПРЕАМБУЛА (ИСПРАВЛЕНИЕ ОШИБОК И ВЫРАВНИВАНИЕ ТЕКСТА)
% ==============================================================================

\usepackage{enumitem}
\usepackage{microtype}

\tolerance=2000
\emergencystretch=1.5em
\hyphenpenalty=500
\hbadness=10000

\microtypesetup{
    expansion=true,
    protrusion=true
}

\setlist{
    nosep,
    itemsep=0.5ex,
    parsep=0.5ex,
    topsep=0.5ex,
    partopsep=0pt,
    leftmargin=*}

\lccode`\-=`\-
\def\hyph{\penalty10000\hskip0pt\penalty8000-\penalty10000\hskip0pt}

% --- Настройка межстрочного интервала ---
\usepackage{setspace}
\onehalfspacing

% --- Геометрия страницы ---
\usepackage[
    left=10mm,
    right=20mm,
    top=20mm,
    bottom=20mm,
    headsep=10mm,
    footskip=10mm
]{geometry}

% --- Мощная математика (AMS + Mathtools) ---
\usepackage{amsmath, amsfonts, amssymb, amscd}
\usepackage{mathtools}
\usepackage{mathrsfs}
\usepackage{bm}

% --- Специфично для матричных вычислений ---
\usepackage{nicematrix}
\usepackage{blkarray}

% --- Алгоритмы (Псевдокод) ---
\usepackage[ruled,vlined,linesnumbered]{algorithm2e}
\SetKwInput{KwData}{Вход}
\SetKwInput{KwResult}{Выход}

% --- Теоремы, леммы, определения ---
\usepackage{amsthm}
\usepackage{thmtools}

\theoremstyle{plain}
\newtheorem{theorem}{Теорема}[section]
\newtheorem{lemma}[theorem]{Лемма}
\newtheorem{corollary}[theorem]{Следствие}
\newtheorem{proposition}[theorem]{Утверждение}

\theoremstyle{definition}
\newtheorem{definition}{Определение}[section]
\newtheorem{example}{Пример}[section]
\newtheorem{problem}{Задача}[section]

\theoremstyle{remark}
\newtheorem{remark}{Замечание}[section]

% --- Графика и таблицы ---
\usepackage{graphicx}
\usepackage{xcolor}
\usepackage{booktabs}
\usepackage{caption}
\usepackage{subcaption}
\usepackage{tikz}
\usetikzlibrary{calc, matrix, arrows.meta, positioning}

% --- Списки ---
\usepackage{enumitem}
\setlist{nosep, itemsep=1ex, leftmargin=*}

% --- Ссылки и навигация ---
\usepackage[colorlinks=true, linkcolor=blue, citecolor=magenta, urlcolor=teal]{hyperref}
\usepackage[noabbrev, russian]{cleveref}

% ==============================================================================
% ЦВЕТОВАЯ ПАЛИТРА ДОКУМЕНТА
% ==============================================================================

\definecolor{TCBBlue}{RGB}{30, 80, 150}
\definecolor{TCBBlueBg}{RGB}{235, 242, 252}
\definecolor{TCBBlueLine}{RGB}{120, 170, 230}

\definecolor{TCBGreen}{RGB}{25, 110, 60}
\definecolor{TCBGreenBg}{RGB}{234, 247, 238}
\definecolor{TCBGreenLine}{RGB}{100, 190, 130}

\definecolor{TCBPurple}{RGB}{90, 40, 140}
\definecolor{TCBPurpleBg}{RGB}{244, 238, 252}
\definecolor{TCBPurpleLine}{RGB}{180, 140, 220}

\definecolor{TCBTeal}{RGB}{15, 105, 120}
\definecolor{TCBTealBg}{RGB}{232, 247, 249}

\definecolor{TCBOrange}{RGB}{180, 80, 20}
\definecolor{TCBOrangeBg}{RGB}{253, 243, 232}

\definecolor{TCBGray}{RGB}{55, 60, 70}
\definecolor{TCBGrayBg}{RGB}{245, 245, 247}

\definecolor{TCBGold}{RGB}{155, 115, 0}
\definecolor{TCBGoldBg}{RGB}{254, 250, 235}

\definecolor{TCBSlate}{RGB}{50, 55, 75}       % Тёмно-сланцевый — доказательства
\definecolor{TCBSlateBg}{RGB}{240, 241, 245}  % Холодный светлый фон

% ==============================================================================
% TCOLORBOX: ЗАГРУЗКА И БАЗОВЫЙ СТИЛЬ
% ==============================================================================
% АРХИТЕКТУРНОЕ РЕШЕНИЕ: все рамки — строго прямоугольники (sharp corners)
% с левой цветной полосой. Ни одно окружение НЕ использует arc или скруглённые углы.
% Это обеспечивает единый академический стиль документа.
% ==============================================================================

\usepackage[most]{tcolorbox}

% ------------------------------------------------------------------------------
% Единственный базовый стиль: прямоугольник + левая цветная полоса.
% ВСЕ окружения документа строятся на этом стиле.
% #1 — цвет полосы, #2 — цвет фона, #3 — цвет заголовка
% ------------------------------------------------------------------------------
\tcbset{
    tcb@leftbar/.style n args={3}{
        enhanced,
        breakable,
        sharp corners,          % ОБЯЗАТЕЛЬНО: все углы прямые, никаких arc=...
        boxrule=0pt,
        frame hidden,
        borderline west={4pt}{0pt}{#1},
        colback=#2,
        colframe=#1,
        coltitle=#3,
        fonttitle=\bfseries\sffamily,
        % --- ИСПРАВЛЕНИЕ: Встроенный заголовок ---
        colbacktitle=#2,   % Фон заголовка совпадает с фоном блока
        titlerule=0pt,     % Убираем линию между заголовком и текстом
        toptitle=2pt,      % Внутренний отступ сверху до заголовка
        bottomtitle=1pt,   % Внутренний отступ от заголовка до текста
        % -----------------------------------------
        top=4pt, bottom=6pt, left=8pt, right=6pt,
        before skip=10pt, after skip=10pt,
    }
}

% ==============================================================================
% TCOLORBOX: АКАДЕМИЧЕСКИЕ ОКРУЖЕНИЯ
% ==============================================================================

% --- 1. ОПРЕДЕЛЕНИЕ (левая синяя полоса) ---
\newtcolorbox[
    auto counter,
    number within=section,
    crefname={определение}{определения},
    Crefname={Определение}{Определения},
]{tcbdefinition}[1][]{
    tcb@leftbar={TCBBlue}{TCBBlueBg}{TCBBlue},
    title={Определение~\thetcbcounter\ifx\relax#1\relax\else\ (\textit{#1})\fi},
    label type=tcbdefinition,
}

% --- 2. ТЕОРЕМА (левая зелёная полоса) ---
\newtcolorbox[
    auto counter,
    number within=section,
    crefname={теорема}{теоремы},
    Crefname={Теорема}{Теоремы},
]{tcbtheorem}[1][]{
    tcb@leftbar={TCBGreen}{TCBGreenBg}{TCBGreen},
    title={Теорема~\thetcbcounter\ifx\relax#1\relax\else\ (\textit{#1})\fi},
    label type=tcbtheorem,
}

% --- 3. ЛЕММА (левая пурпурная полоса) ---
\newtcolorbox[
    auto counter,
    number within=section,
    crefname={лемма}{леммы},
    Crefname={Лемма}{Леммы},
]{tcblemma}[1][]{
    tcb@leftbar={TCBPurple}{TCBPurpleBg}{TCBPurple},
    title={Лемма~\thetcbcounter\ifx\relax#1\relax\else\ (\textit{#1})\fi},
    label type=tcblemma,
}

% --- 4. СЛЕДСТВИЕ (левая бирюзовая полоса) ---
\newtcolorbox[
    auto counter,
    number within=section,
    crefname={следствие}{следствия},
    Crefname={Следствие}{Следствия},
]{tcbcorollary}[1][]{
    tcb@leftbar={TCBTeal}{TCBTealBg}{TCBTeal},
    title={Следствие~\thetcbcounter\ifx\relax#1\relax\else\ (\textit{#1})\fi},
    label type=tcbcorollary,
}

% --- 5. ДОКАЗАТЕЛЬСТВО (левая сланцевая полоса, знак QED справа) ---
% Использование: \begin{tcbproof} ... \end{tcbproof}
% Знак \(\square\) добавляется АВТОМАТИЧЕСКИ — вручную \qed НЕ ставить.
\newtcolorbox{tcbproof}[1][]{
    tcb@leftbar={TCBSlate}{TCBSlateBg}{TCBSlate},
    title={Доказательство\ifx\relax#1\relax\else\ (\textit{#1})\fi},
    after upper={\hfill\(\square\)},
}

% --- 6. ЗАМЕЧАНИЕ (левая серая полоса) ---
\newtcolorbox{tcbremark}[1][]{
    tcb@leftbar={TCBGray!60}{TCBGrayBg}{TCBGray},
    title={Замечание\ifx\relax#1\relax\else\ (\textit{#1})\fi},
}

% --- 7. ПРЕДУПРЕЖДЕНИЕ / ВНИМАНИЕ (левая оранжевая полоса) ---
\newtcolorbox{tcbwarning}[1][]{
    tcb@leftbar={TCBOrange}{TCBOrangeBg}{TCBOrange},
    title={\(\vartriangle\)\ Внимание\ifx\relax#1\relax\else:\ \textit{#1}\fi},
}

% --- 8. ПРИМЕР (левая серая полоса, белый фон) ---
\newtcolorbox[
    auto counter,
    number within=section,
    crefname={пример}{примера},
    Crefname={Пример}{Примера},
]{tcbexample}[1][]{
    tcb@leftbar={TCBGray!50}{white}{TCBGray},
    title={Пример~\thetcbcounter\ifx\relax#1\relax\else\ (\textit{#1})\fi},
    label type=tcbexample,
}

% --- 9. КЛЮЧЕВАЯ ИДЕЯ / ИНСАЙТ (левая золотая полоса) ---
\newtcolorbox{tcbinsight}[1][]{
    tcb@leftbar={TCBGold}{TCBGoldBg}{TCBGold},
    title={\(\star\)\ Ключевая идея\ifx\relax#1\relax\else:\ \textit{#1}\fi},
}

% --- 10. АЛГОРИТМ (левая синяя полоса) ---
% [КРИТИЧЕСКИ ВАЖНО] Использовать для ЛЮБОГО алгоритма, псевдокода,
% пошаговой инструкции и нумерованного списка последовательных требований/шагов.
% Использование: \begin{tcbalgorithmbox}[Название алгоритма] ... \end{tcbalgorithmbox}
\newtcolorbox[
    auto counter,
    number within=section,
    crefname={алгоритм}{алгоритма},
    Crefname={Алгоритм}{Алгоритма},
]{tcbalgorithmbox}[1][]{
    tcb@leftbar={TCBBlue}{TCBBlueBg}{TCBBlue},
    title={Алгоритм~\thetcbcounter\ifx\relax#1\relax\else:\ \textit{#1}\fi},
    fontupper=\small,
    label type=tcbalgorithmbox,
}

% ==============================================================================
% КАСТОМНЫЕ МАКРОСЫ И ОПЕРАТОРЫ ДЛЯ МАТРИЧНЫХ ВЫЧИСЛЕНИЙ
% ==============================================================================

\newcommand{\R}{\mathbb{R}}
\renewcommand{\C}{\mathbb{C}}
\newcommand{\N}{\mathbb{N}}
\newcommand{\Z}{\mathbb{Z}}

\DeclareMathOperator{\rank}{rank}
\DeclareMathOperator{\tr}{tr}
\DeclareMathOperator{\diag}{diag}
\DeclareMathOperator{\im}{Im}
\undef\ker
\DeclareMathOperator{\ker}{Ker}
\DeclareMathOperator{\Span}{span}
\DeclareMathOperator{\cond}{cond}
\DeclareMathOperator{\sign}{sign}

\newcommand{\T}{^{\mathsf{T}}}
\newcommand{\Htrans}{^{\mathsf{H}}}
\let\H\Htrans

\DeclarePairedDelimiter{\abs}{\lvert}{\rvert}
\DeclarePairedDelimiter{\norm}{\lVert}{\rVert}
\DeclarePairedDelimiter{\inner}{\langle}{\rangle}
\DeclarePairedDelimiter{\ceil}{\lceil}{\rceil}
\DeclarePairedDelimiter{\floor}{\lfloor}{\rfloor}

\makeatletter
\let\oldabs\abs
\def\abs{\@ifstar{\oldabs}{\oldabs*}}
\let\oldnorm\norm
\def\norm{\@ifstar{\oldnorm}{\oldnorm*}}
\makeatother

\usepackage{import}

% ==============================================================================
% КОНЕЦ ПРЕАМБУЛЫ. НАЧАЛО ДОКУМЕНТА
% ==============================================================================

\title{\textbf{Основы матричных вычислений}\\ \Large Семинар / Лекция}
\author{Название университета / Факультет}
\date{\today}

\begin{document}
\tableofcontents
\newpage
\maketitle

% ==============================================================================
% lec9.tex — Демонстрационный файл: все правила вёрстки + матричные примеры
% Каждый блок явно помечен, какое правило промпта он иллюстрирует.
% ==============================================================================

\begin{abstract}
    \noindent Данный документ демонстрирует корректное использование преамбулы.
    Все генерации ИИ обязаны следовать строгой семантической структуре,
    показанной ниже.
\end{abstract}

\vspace{5mm}

% ==============================================================================
\section{Сингулярное разложение (SVD)}
\label{sec:svd}
% ==============================================================================

% ------------------------------------------------------------------------------
% ПРАВИЛО 1: никакого Markdown — только LaTeX-команды.
% ПРАВИЛО 5: tcbdefinition для определений.
% ПРАВИЛО 8: метка def:svd по соглашению def:...
% ------------------------------------------------------------------------------
\begin{tcbdefinition}[Сингулярное разложение]
\label{def:svd}
    Пусть матрица $\bm{A} \in \R^{m \times n}$.
    \textit{Сингулярным разложением} (SVD) матрицы $\bm{A}$ называется
    факторизация вида
    % ПРАВИЛО 3: equation с меткой, никаких $$...$$
    \begin{equation}
        \bm{A} = \bm{U}\,\bm{\Sigma}\,\bm{V}\T,
        \label{eq:svd}
    \end{equation}
    где $\bm{U} \in \R^{m \times m}$ и $\bm{V} \in \R^{n \times n}$ —
    ортогональные матрицы, а $\bm{\Sigma} \in \R^{m \times n}$ —
    псевдодиагональная матрица с сингулярными числами
    $\sigma_1 \geq \sigma_2 \geq \cdots \geq \sigma_r > 0$,
    $r = \rank(\bm{A})$.
\end{tcbdefinition}

% ПРАВИЛО 8: ссылки ТОЛЬКО через \cref / \Cref, никогда через голый \ref
Согласно \cref{def:svd}, разложение~\cref{eq:svd} существует для любой
вещественной матрицы.
\Cref{thm:svd-exist} ниже доказывает это утверждение.

% ПРАВИЛО 3: матрица через pNiceMatrix с подписями строк/столбцов
% ПРАВИЛО 4: \vspace для отступов, не пустые строки
Матрица $\bm{\Sigma}$ из~\cref{eq:svd} имеет вид
\begin{equation}
    \bm{\Sigma} =
    \begin{pNiceMatrix}[first-row, first-col, margin]
              & 1          & \Cdots  &              & n      \\
        1     & \sigma_1   &         &              &        \\
        \Vdots&            & \Ddots  &              & \Vdots \\
              &            &         & \sigma_r     &        \\
        m     &            & \Cdots  &              & 0
    \end{pNiceMatrix},
    \label{eq:sigma}
\end{equation}
где нули занимают все позиции вне главной диагонали.

% ------------------------------------------------------------------------------
% ПРАВИЛО 5: tcbtheorem для утверждений.
% ПРАВИЛО 8: метка thm:svd-exist по соглашению thm:...
% ------------------------------------------------------------------------------
\begin{tcbtheorem}[Существование SVD]
\label{thm:svd-exist}
    Для любой матрицы $\bm{A} \in \C^{m \times n}$ существует сингулярное
    разложение $\bm{A} = \bm{U}\,\bm{\Sigma}\,\bm{V}\H$.
    Сингулярные числа $\sigma_i$ определены однозначно;
    при их попарной различности матрицы $\bm{U}$ и $\bm{V}$ единственны
    с точностью до знаков столбцов.
\end{tcbtheorem}

% ------------------------------------------------------------------------------
% ПРАВИЛО 7: доказательство ТОЛЬКО через tcbproof.
%            Знак \square добавляется АВТОМАТИЧЕСКИ — вручную \qed не ставить.
% ------------------------------------------------------------------------------
\begin{tcbproof}[теоремы~\ref{thm:svd-exist}]
    Рассмотрим симметричную матрицу $\bm{A}\T\bm{A} \in \R^{n \times n}$.
    Она положительно полуопределена, поэтому имеет спектральное разложение
    \begin{equation}
        \bm{A}\T\bm{A}
        = \bm{V}\,\diag(\sigma_1^2, \ldots, \sigma_r^2, 0, \ldots, 0)\,\bm{V}\T.
        \label{eq:ata}
    \end{equation}
    Полагая $\bm{u}_i = \bm{A}\bm{v}_i / \sigma_i$ для $i = 1, \ldots, r$
    и дополняя до ортонормированного базиса, получаем матрицу $\bm{U}$.
    Непосредственная проверка показывает, что $\bm{A} = \bm{U}\bm{\Sigma}\bm{V}\T$.
\end{tcbproof}

% ------------------------------------------------------------------------------
% ПРАВИЛО 5: tcbinsight для ключевых идей.
% ------------------------------------------------------------------------------
\begin{tcbinsight}[Геометрический смысл SVD]
    SVD представляет линейное отображение $\bm{A}$ как последовательность
    трёх геометрических преобразований: поворота $\bm{V}\T$, масштабирования
    $\bm{\Sigma}$ и второго поворота $\bm{U}$.
    Сингулярные числа $\sigma_i$ — коэффициенты растяжения вдоль главных осей.
\end{tcbinsight}

% ==============================================================================
\subsection{Свойства SVD и вычислительные аспекты}
\label{sec:svd-props}
% ==============================================================================

% ПРАВИЛО 2: список через \begin{itemize}, никаких дефисов текстом.
% ПРАВИЛО 9: макросы \rank, \norm, \cond, \tr.
Основные свойства сингулярного разложения:
\begin{itemize}
    \item Ранг матрицы равен числу ненулевых сингулярных чисел:
          $\rank(\bm{A}) = r$.
    \item Евклидова норма совпадает с максимальным сингулярным числом:
          $\norm{\bm{A}}_2 = \sigma_1$.
    \item Число обусловленности: $\cond_2(\bm{A}) = \sigma_1 / \sigma_r$.
    \item Фробениусова норма: $\norm{\bm{A}}_F^2 = \tr(\bm{A}\T\bm{A})
          = \sum_{i=1}^r \sigma_i^2$.
\end{itemize}

% ПРАВИЛО 5: tcbremark для пояснений.
\begin{tcbremark}
    Для иллюстрации блочных вычислений ниже используется окружение
    \texttt{pNiceMatrix}. Пакет \texttt{nicematrix} идеально подходит
    для работы с блочными матрицами, пунктирными линиями и сложными
    выделениями.
\end{tcbremark}

% ПРАВИЛО 3: матрица через pNiceMatrix, через \[ ... \], не $$
\[
    \bm{M} =
    \begin{pNiceMatrix}
        a_{11} & a_{12} \\
        a_{21} & a_{22}
    \end{pNiceMatrix}.
\]

% ------------------------------------------------------------------------------
% ПРАВИЛО 6: ЛЮБОЙ алгоритм / пошаговая процедура — ТОЛЬКО в tcbalgorithmbox.
%            Голый enumerate без обёртки ЗАПРЕЩЁН для последовательных шагов.
% ПРАВИЛО 8: метка alg:power-svd по соглашению alg:...
% ------------------------------------------------------------------------------
\begin{tcbalgorithmbox}[Степенная итерация для нахождения $\sigma_1$]
\label{alg:power-svd}
    \textbf{Вход:} матрица $\bm{A} \in \R^{m \times n}$,
    начальный вектор $\bm{x}_0 \in \R^n$, точность $\varepsilon > 0$.\\
    \textbf{Выход:} наибольшее сингулярное число $\sigma_1$
    и правый сингулярный вектор $\bm{v}_1$.
    \begin{enumerate}
        \item Нормировать: $\bm{v} \gets \bm{x}_0 / \norm{\bm{x}_0}$.
        \item Повторять до сходимости:
        \begin{enumerate}
            \item $\bm{u} \gets \bm{A}\bm{v}$;\quad
                  $\sigma \gets \norm{\bm{u}}$;\quad
                  $\bm{u} \gets \bm{u}/\sigma$.
            \item $\bm{v} \gets \bm{A}\T\bm{u}$;\quad
                  $\bm{v} \gets \bm{v}/\norm{\bm{v}}$.
            \item Если $\abs{\sigma_{\mathrm{new}} - \sigma_{\mathrm{old}}}
                  < \varepsilon$ — остановиться.
        \end{enumerate}
        \item Вернуть $(\sigma,\,\bm{v})$.
    \end{enumerate}
\end{tcbalgorithmbox}

% ПРАВИЛО 8: ссылка на алгоритм через \cref
Сходимость \cref{alg:power-svd} определяется отношением $\sigma_2/\sigma_1$:
чем оно меньше, тем быстрее итерация.

% ==============================================================================
\section{Структурные матрицы с корректной сеткой}
\label{sec:structured}
% ==============================================================================

В этом разделе представлены матрицы, где структура блоков строго соответствует
количеству строк и столбцов.

% ==============================================================================
\subsection{Блочно-ленточная матрица}
\label{sec:block-band}
% ==============================================================================

Для корректной работы \texttt{\textbackslash Block\{2-2\}} в нижней части
матрицы явно добавляем пустую строку, расширяя сетку до нужного размера.

\begin{equation}
    \bm{\mathcal{M}} =
    \begin{pNiceMatrix}[margin, columns-width=10mm]
        \Block[draw]{2-2}{\bm{B}_1} & & \bm{0}   & \Cdots & \bm{0}   \\
        & & \Vdots   &        & \Vdots   \\
        \bm{0} & \Cdots & \bm{B}_i & \Cdots & \bm{0}   \\
        \Vdots &        & \Vdots   & \Block[draw]{2-2}{\bm{B}_n} & \\
        \bm{0} & \Cdots & \bm{0}   &        &
    \end{pNiceMatrix}.
    \label{eq:block-band}
\end{equation}

% ==============================================================================
\subsection{Разреженная трёхдиагональная матрица}
\label{sec:tridiag}
% ==============================================================================

\begin{equation}
    \bm{A} =
    \begin{pNiceMatrix}
        a_{11} & a_{12} & 0         & \Cdots    & 0            \\
        a_{21} & a_{22} & a_{23}    &           & \Vdots       \\
        0      & a_{32} & \Ddots    & \Ddots    & 0            \\
        \Vdots &        & \Ddots    & \Ddots    & a_{n-1,n}    \\
        0      & \Cdots & 0         & a_{n,n-1} & a_{nn}
    \end{pNiceMatrix}.
    \label{eq:tridiag}
\end{equation}

% ==============================================================================
\subsection{Блочная матрица с дополнением Шура}
\label{sec:schur}
% ==============================================================================

Тонкие внутренние линии визуально разделяют подматрицы.

\begin{equation}
    \bm{K} =
    \begin{pNiceMatrix}[
        margin,
        hvlines-except-borders,
        rules/color=gray!30,
        columns-width=12mm
    ]
        \Block{2-2}{\bm{A}_{11}} & & \Block{2-2}{\bm{A}_{12}} & \\
        & & & \\
        \Block{2-2}{\bm{A}_{21}} & & \Block[fill=red!5]{2-2}{\bm{S}} & \\
        & & &
    \end{pNiceMatrix}.
    \label{eq:schur-block}
\end{equation}

% ПРАВИЛО 5: tcbcorollary.
% ПРАВИЛО 8: метка cor:schur по соглашению cor:...
\begin{tcbcorollary}[Дополнение Шура]
\label{cor:schur}
    Если матрица $\bm{A}_{11}$ невырождена, то дополнение Шура блока
    $\bm{A}_{11}$ в матрице~\cref{eq:schur-block} равно
    \[
        \bm{S} = \bm{A}_{22} - \bm{A}_{21}\,\bm{A}_{11}^{-1}\,\bm{A}_{12}.
    \]
\end{tcbcorollary}

% ПРАВИЛО 7: доказательство следствия — tcbproof.
\begin{tcbproof}[следствия~\ref{cor:schur}]
    Блочное разложение
    \begin{equation}
        \bm{K} =
        \begin{pNiceMatrix}
            \bm{I} & \bm{0} \\
            \bm{A}_{21}\bm{A}_{11}^{-1} & \bm{I}
        \end{pNiceMatrix}
        \begin{pNiceMatrix}
            \bm{A}_{11} & \bm{0} \\
            \bm{0}      & \bm{S}
        \end{pNiceMatrix}
        \begin{pNiceMatrix}
            \bm{I} & \bm{A}_{11}^{-1}\bm{A}_{12} \\
            \bm{0} & \bm{I}
        \end{pNiceMatrix}
        \label{eq:schur-proof}
    \end{equation}
    проверяется прямым умножением.
    Из~\cref{eq:schur-proof} немедленно следует
    $\det(\bm{K}) = \det(\bm{A}_{11})\cdot\det(\bm{S})$.
\end{tcbproof}

% ==============================================================================
\subsection{Двухдиагональная матрица Голуба–Кахана}
\label{sec:bidiag}
% ==============================================================================

Применяется на первом этапе алгоритмов вычисления SVD.

\begin{equation}
    \bm{B} =
    \begin{pNiceMatrix}
        \alpha_1 & \beta_1  & 0        & \Cdots & 0           \\
        0        & \alpha_2 & \beta_2   &        & \Vdots      \\
        \Vdots   & \Ddots   & \Ddots    & \Ddots & 0           \\
        \Vdots   &          & \Ddots    & \alpha_{n-1} & \beta_{n-1} \\
        0        & \Cdots   & \Cdots    & 0      & \alpha_n
    \end{pNiceMatrix}.
    \label{eq:bidiag}
\end{equation}

% ПРАВИЛО 5: tcbwarning — для предупреждений, ОБЯЗАТЕЛЬНО в блоке.
\begin{tcbwarning}[Типичная ошибка: Block too large]
    При возникновении ошибки компилятора \texttt{Block too large} проверить:
    \begin{enumerate}
        \item Соответствует ли число символов \texttt{\&} размеру блока.
              Для блока \texttt{\{2-2\}} в последнем столбце нужно
              добавить ещё один пустой столбец справа.
        \item Добавлена ли пустая строка \texttt{\textbackslash\textbackslash}
              после блока, если он расположен в самом низу матрицы.
    \end{enumerate}
\end{tcbwarning}

% ==============================================================================
\section{Продвинутые примеры с пакетом \texttt{nicematrix}}
\label{sec:nicematrix}
% ==============================================================================

\begin{tcbremark}
    Пакет \texttt{nicematrix} идеально подходит для работы с блочными
    матрицами, пунктирными линиями и сложными выделениями.
    Ниже представлены структурные примеры, расширяющие базовый функционал.
\end{tcbremark}

% ==============================================================================

\begin{equation}
    \bm{D} =
    \begin{pNiceMatrix}[margin]
        % \Block с пустыми скобками {} красит фон, но оставляет текст ячеек независимым
        \Block[fill=blue!10]{2-2}{} d_{11} & d_{12} & 0      & 0      & \Cdots & 0      \\
        d_{21} & d_{22} & 0      & 0      &        & \Vdots \\
        0      & 0      & \Block[fill=green!10]{2-2}{} d_{33} & d_{34} & \Ddots &        \\
        0      & 0      & d_{43} & d_{44} & \Ddots & 0      \\
        \Vdots &        & \Ddots & \Ddots & d_{55} & d_{56} \\
        0      & \Cdots &        & 0      & d_{65} & d_{66}
    \end{pNiceMatrix}.
    \label{eq:block-diag}
\end{equation}

% ==============================================================================
\subsection{Расширенная матрица для системы линейных уравнений}
\label{sec:augmented}
% ==============================================================================

Окружение \texttt{pNiceArray} позволяет задавать спецификацию столбцов
(как в стандартном \texttt{array}) — для вертикальной черты перед столбцом
свободных членов.

\begin{equation}
    \bm{M}_{\mathrm{aug}} =
    \begin{pNiceArray}{ccc|c}[margin]
        a_{11} & a_{12} & a_{13} & b_1 \\
        a_{21} & a_{22} & a_{23} & b_2 \\
        a_{31} & a_{32} & a_{33} & b_3
    \end{pNiceArray}.
    \label{eq:augmented}
\end{equation}

% ==============================================================================
\subsection{Матрица переходных вероятностей с подписями состояний}
\label{sec:markov}
% ==============================================================================

Опции \texttt{first-row} и \texttt{first-col} выносят первую строку
и первый столбец \emph{за} пределы математических скобок матрицы.

\begin{equation}
    \bm{P} =
    \begin{pNiceMatrix}[first-row, first-col, margin]
             & x_1    & x_2    & x_3    \\
        y_1  & p_{11} & p_{12} & p_{13} \\
        y_2  & p_{21} & p_{22} & p_{23} \\
        y_3  & p_{31} & p_{32} & p_{33}
    \end{pNiceMatrix}.
    \label{eq:markov}
\end{equation}

% ПРАВИЛО 5: tcblemma.
% ПРАВИЛО 8: метка lem:stochastic по соглашению lem:...
\begin{tcblemma}[Стохастичность матрицы переходов]
\label{lem:stochastic}
    Матрица $\bm{P}$ из~\cref{eq:markov} является стохастической тогда
    и только тогда, когда
    \[
        p_{ij} \geq 0 \quad \text{и} \quad \sum_{j} p_{ij} = 1
        \quad \forall\, i.
    \]
\end{tcblemma}

\begin{tcbproof}[леммы~\ref{lem:stochastic}]
    \textit{Необходимость.} Если $\bm{P}$ стохастична, то по определению
    все элементы неотрицательны, а сумма элементов каждой строки равна
    единице, что в точности совпадает с формулой $\sum_j p_{ij} = 1$.\\
    \textit{Достаточность.} При выполнении указанных условий вектор
    $\mathbf{1}$ является правым собственным вектором $\bm{P}$ с
    собственным значением $1$, откуда следует сохранение вероятностной
    меры при умножении.
\end{tcbproof}

% ==============================================================================
\subsection{Вложенные блоки с явной отрисовкой границ}
\label{sec:nested-blocks}
% ==============================================================================

Опция \texttt{draw} обводит блок рамкой; \texttt{fill} задаёт фоновый цвет.

\begin{equation}
    \bm{H} =
    \begin{pNiceMatrix}[margin, columns-width=12mm]
        \Block[draw, fill=blue!5]{2-2}{\bm{H}_{11}} & &
        \Block[draw]{2-2}{\bm{H}_{12}} & \\
        & & & \\
        \Block[draw]{2-2}{\bm{H}_{21}} & &
        \Block[draw, fill=red!5]{2-2}{\bm{H}_{22}} & \\
        & & &
    \end{pNiceMatrix}.
    \label{eq:nested-blocks}
\end{equation}

% ==============================================================================
\subsection{Верхнетреугольная матрица QR-разложения}
\label{sec:qr}
% ==============================================================================

\texttt{\textbackslash CodeAfter} позволяет рисовать поверх уже
расставленных элементов; пунктирная диагональная линия обозначает границу
нулевой нижней части.

\begin{equation}
    \bm{R} =
    \begin{pNiceMatrix}[margin]
        r_{11} & r_{12} & r_{13} & \Cdots & r_{1n}      \\
        0      & r_{22} & r_{23} & \Cdots & r_{2n}      \\
        \Vdots & \Ddots & \ddots &        & \Vdots      \\
               &        &        &        & r_{n-1,\,n} \\
        0      & \Cdots &        & 0      & r_{nn}
    \CodeAfter
        \tikz \draw [dashed, gray]
            (2-|1) -| (3-|2) -| (4-|3) -| (5-|4) -| (6-|5);
    \end{pNiceMatrix}.
    \label{eq:qr-upper}
\end{equation}

% ПРАВИЛО 5: tcbexample.
% ПРАВИЛО 8: метка ex:qr по соглашению ex:...
\begin{tcbexample}[Применение QR-разложения к решению СЛАУ]
\label{ex:qr}
    Для матрицы $\bm{A} \in \R^{3 \times 3}$ процесс Грама–Шмидта даёт
    ортогональный множитель $\bm{Q}$ и верхнетреугольный множитель
    $\bm{R}$ (\cref{eq:qr-upper}) такой, что $\bm{A} = \bm{Q}\bm{R}$.
    Система $\bm{A}\bm{x} = \bm{b}$ сводится к
    $\bm{R}\bm{x} = \bm{Q}\T\bm{b}$, решаемой обратной подстановкой
    за $O(n^2)$ операций.
\end{tcbexample}

% ==============================================================================
\section{Сводная демонстрация правил 6 и 9}
\label{sec:rules-demo}
% ==============================================================================

% ------------------------------------------------------------------------------
% ПРАВИЛО 6: последовательность требований — ТОЛЬКО в tcbalgorithmbox,
%            никогда в голом enumerate вне обёртки.
% ПРАВИЛО 8: метка alg:golub-kahan
% ------------------------------------------------------------------------------
\begin{tcbalgorithmbox}[Алгоритм Голуба–Кахана (бидиагонализация)]
\label{alg:golub-kahan}
    \textbf{Вход:} $\bm{A} \in \R^{m \times n}$, $m \geq n$.\\
    \textbf{Выход:} ортогональные $\bm{U}$, $\bm{V}$;
    бидиагональная матрица $\bm{B}$ вида~\cref{eq:bidiag}.
    \begin{enumerate}
        \item Выбрать начальный вектор $\bm{p}_1$ с $\norm{\bm{p}_1} = 1$.
        \item Для $j = 1, 2, \ldots, n$ выполнить:
        \begin{enumerate}
            \item $\bm{r}_j \gets \bm{A}\T\bm{p}_j
                  - \beta_{j-1}\bm{q}_{j-1}$.
            \item $\alpha_j \gets \norm{\bm{r}_j}$;\quad
                  $\bm{q}_j \gets \bm{r}_j / \alpha_j$.
            \item $\bm{s}_j \gets \bm{A}\bm{q}_j - \alpha_j\bm{p}_j$.
            \item $\beta_j \gets \norm{\bm{s}_j}$;\quad
                  $\bm{p}_{j+1} \gets \bm{s}_j / \beta_j$.
        \end{enumerate}
        \item Собрать $\bm{B}$ из $\{\alpha_j\}$ и $\{\beta_j\}$
              по образцу~\cref{eq:bidiag}.
    \end{enumerate}
\end{tcbalgorithmbox}

% ПРАВИЛО 8: ссылки на несколько объектов через \cref
После выполнения \cref{alg:golub-kahan} к матрице $\bm{B}$ применяется
\cref{alg:power-svd}, что существенно снижает вычислительную стоимость
по сравнению с прямым SVD исходной матрицы $\bm{A}$ (см.~\cref{def:svd}).

% ПРАВИЛО 9: макросы \norm, \inner, \abs, \rank, \tr в тексте
Для матриц из данного раздела используются следующие нормы и операторы:
\begin{itemize}
    \item $\norm{\bm{A}}_2 = \sigma_1(\bm{A})$ — спектральная норма;
    \item $\norm{\bm{A}}_F = \sqrt{\tr(\bm{A}\T\bm{A})}$ — норма Фробениуса;
    \item $\inner{\bm{u}, \bm{v}} = \bm{u}\T\bm{v}$ — стандартное скалярное
          произведение в $\R^n$;
    \item $\abs{\lambda_i(\bm{A})} \leq \norm{\bm{A}}_2$ — оценка
          собственных значений через норму.
\end{itemize}

% ПРАВИЛО 5: tcbinsight
\begin{tcbinsight}[Связь SVD и спектрального разложения]
    Из~\cref{eq:ata} в доказательстве \cref{thm:svd-exist} видно:
    сингулярные числа $\bm{A}$ — это квадратные корни из собственных
    значений $\bm{A}\T\bm{A}$, а столбцы $\bm{V}$ — соответствующие
    собственные векторы.
    Это связывает \cref{def:svd} со спектральной теорией симметричных матриц.
\end{tcbinsight}

\vspace{6pt}

% Справочная таблица (booktabs, без лишних вертикальных линий)
\begin{table}[h]
\centering
\caption{Справочник окружений и их назначений}
\label{tab:envs}
\begin{tabular}{@{}lll@{}}
\toprule
\textbf{Окружение}        & \textbf{Цвет полосы} & \textbf{Назначение}                    \\
\midrule
\texttt{tcbdefinition}    & синий                & Определения понятий                    \\
\texttt{tcbtheorem}       & зелёный              & Утверждения с доказательством          \\
\texttt{tcblemma}         & пурпурный            & Вспомогательные леммы                  \\
\texttt{tcbcorollary}     & бирюзовый            & Следствия из теорем                    \\
\texttt{tcbproof}         & сланцевый            & Доказательства (авто-\(\square\))      \\
\texttt{tcbremark}        & серый                & Замечания и пояснения                  \\
\texttt{tcbwarning}       & оранжевый            & Предупреждения и типичные ловушки      \\
\texttt{tcbexample}       & серый/белый          & Числовые и задачные примеры            \\
\texttt{tcbinsight}       & золотой              & Ключевые идеи                          \\
\texttt{tcbalgorithmbox}  & синий                & Алгоритмы и пошаговые процедуры        \\
\bottomrule
\end{tabular}
\end{table}

\end{document}
